\subsection*{A selecting principle for the bond vertex}
Choose the canonical coordinate CAR field $c$, the proper regulator
$R_\psi=\exp[-(\psi^{-1}D\psi^{-1})^2/(2\Lambda^2)]$,
physical range $\ell$, and a ruler-distance scalar kernel.
For the compensated continuum dilation
$\psi_\lambda(x)=\lambda^{1/2}\psi(\lambda x)$,
$N_\lambda(x)=N(\lambda x)$ and
$c_\lambda(x)=\lambda^{3/2}c(\lambda x)$, with
$T_\lambda f(x)=\lambda^{3/2}f(\lambda x)$,
spectral calculus gives $R_{\psi_\lambda}=T_\lambda R_\psi T_\lambda^\dagger$.
The distance satisfies
$d_{\psi_\lambda}(x,y)=d_\psi(\lambda x,\lambda y)$.
Thus the frame bilinear
$j_{B,a}=(R_\psi c)^\dagger\alpha_a(R_\psi c)$ has weight 3,
and its bilocal interaction with contraction $\delta_{ab}$ has
weight $3+3-6=0$.

This convention follows from a more explicit, sufficient microscopic
principle: use the proper-volume field
$\chi_\Lambda=\psi^{-3}R_\psi c$, proper measures
$dV=\psi^6d^3x$, and contract the 2 frame bilinears with the
chosen common frame alignment.
The 2 volume factors cancel their 2 density factors exactly,
leaving the coordinate interaction built from $j_B$.
This principle specifies the vertex; it is not a consequence of H2.
In particular $j_B$ is not identified with the full Noether current
or momentum density of the completed nonlocal law.

For a counterfamily replace each endpoint by
$j_b=\psi^{-b}j_B$, retaining the same scalar kernel and contraction.
The seed weight becomes $-b$.
An independent infinitesimal check gives the pair-field weight
$6-b$; integration of
$(x\cdot\nabla_x+y\cdot\nabla_y)(\mathcal K j_bj_b)$
contributes $-6$, again leaving $-b$.
[X: finite field/kernel pullback; canonical pair-kernel virial.]
Each $b$ permits the same positive completion and retains vacuum
and 1-excitation protection, lapse affinity and Gram clock sources.
These properties therefore do not select $b=0$.
The apparent alternative $b=2$ is exactly the original law if its
contraction is also changed:
\[
 \begin{aligned}
 \mathcal G_{ij'}(x,y)&=\psi^2(x)\psi^2(y)\delta_{ij'},
 &j_{\rm tr}^i&=\psi^{-2}j_{B,i},\\
 \mathcal G_{ij'}j_{\rm tr}^i(x)j_{\rm tr}^{j'}(y)
   &=\delta_{ab}j_{B,a}(x)j_{B,b}(y).
 \end{aligned}
\]
Changing only the current is a different interaction.

For the transported spectral filter and the physical unit-lapse
positive-sum denominator, the full completion inherits this covariance:
\[
 \mathscr U_\lambda^\dagger
 H_{\rm exc}^{(b)}[N_\lambda,\psi_\lambda;g_B]\mathscr U_\lambda
   =H_{\rm exc}^{(b)}[N,\psi;\lambda^{-b}g_B].
\]
Indeed the pair-energy Sylvester equation
$E_2S_2+S_2E_2=V_{\rm ret}$ intertwines; its unique solution and
the Hermitian similarity exponential intertwine with it.
All other physical parameters are fixed.
Independently, differentiating the Sylvester equation gives
$\delta S=[G,S]-b g_B\partial_{g_B}S$ for the anti-Hermitian
canonical dilation generator $G$; the exponential Fr\'echet derivative gives the same
formula for $H_{\rm exc}$.
For an isolated differentiable stationary state with vanishing
matter dilation boundary terms, set
$F=\langle H_{\rm exc}^{(b)}\rangle$ and take its derivatives at
fixed canonical state.  The complete source integral is
\[
 \begin{aligned}
 \mathcal D_m^{(b)}
 &=\int [F_N\,x\cdot\nabla N+
            F_\psi(\psi/2+x\cdot\nabla\psi)]\,d^3x\\
 &=-b g_B\langle\partial_{g_B}H_{\rm exc}^{(b)}\rangle .
 \end{aligned}
\]
[X: finite spectral transport; complete Sylvester/Fr\'echet response.]
It vanishes identically for the selected $b=0$ construction.
At finite coupling the right side is not generally
$-b\langle V_B\rangle$.
The regulator, moving free projector, excitation filter and dressing
derivatives are part of this identity.
It is a continuum matter identity before gravitational rows or masses
are used, not a fixed-spacing symmetry or a physical mass-lock theorem.

\paragraph*{SOURCES.}
\begingroup\footnotesize
\path{<internal>/P15_current_and_dilation.md},
sections 1--5, actual attachment read completely: regulator covariance,
endpoint coframes, the $b$ counterfamily, physical unit-lapse
denominator, moving-filter term and all-coupling defect.
\path{<internal>/rederivation_final3_ledger.tex}
and
\path{<internal>/redteam_final3_ledger.tex},
conditional Ward derivations: finite spectral intertwining,
canonical generator/virial, and complete Fr\'echet response.
These named methods support the completed $b=0$ identity [X];
the $b$ extension changes its seed covariance as displayed.
\path{<internal>/STATE.tex},
matter-law and dilation entries: the source-only dilation-weight
statement is sharpened here to the full coupling derivative.
Positive rulers, a common differentiable spectral domain and a valid
stationary-state commutator trace are required.
No source-reported numerical checks are counted as executions here.
\endgroup
% Governing source: <internal>/GOAL_alt.md,
% duty 4(3). The chosen scalar kernel may include the stated symmetric
% clock placement; N has zero compensated weight.
% Direct virial check uses the canonical field variation
% delta c=(3/2+x.grad)c, obtained from its dilation generator.
% Two endpoints give 6 before the explicit psi^(-b) factors give -b.
% With fixed masses, ell, Lambda and g, delta R=[G,R] follows also
% from the Duhamel derivative of exp[-K^2/(2Lambda^2)].
% The complete moving-filter variation is
% delta Pi=[ad_G,Pi], rather than zero.
% E2=epsilon tensor I+I tensor epsilon on the antisymmetric pair
% space, epsilon the physical B[1,psi] excitation operator.
% The four-leg denominator is its positive Sylvester sum, not |A|.
% At fixed finite excitation number, bounded retained pair operators
% and a positive one-body gap suffice for the exponential construction;
% no bounded full-Fock similarity is assumed.
